% Experimental Setup section
% Numbers sourced from: research/memory/findings.md (F1–F25), research/specs/000[3-14].md (pre-registered gates),
% results/exp0[1-14]_*/ (seed-level results), and docs/decisions/000[3-7].md (protocol amendments).
% Hardware: GB10 (DGX Spark, 121GB unified memory, thermal governor ~91°C defect) + RTX 5090 (32GB) for training.
% Evaluation: BLiMP (1000-pair subset of minimal pairs) at 11M–121M; lm-evaluation-harness at 1–3B (Qwen3-1.7B base rates).

\section{Experimental Setup and Closure Protocol}
\label{sec:setup}

This paper adopts a pre-registered empirical program with adversarial audit, as detailed in Section~\ref{sec:method}.
Every hypothesis was registered with a quantitative, pass/fail gate in a specification document version-controlled
before its experiment ran. The closure contract mandates six layers: (i) all experiment arms complete; (ii) no
verdict based on fewer than three md5-distinct random seeds; (iii) every reported number traced to a results file
carrying a resolved-config SHA-256 hash and git commit; (iv) a resolved adversarial review that recomputes all
results from raw data and audits like-for-like fairness; (v) all findings and review verdicts recorded in the
findings ledger; and (vi) operator sign-off before merge to main. Figure~\ref{fig:protocol} summarizes the workflow:
an operator registers a research question and its quantitative gate; the experimenter designs a specification (typically
a YAML config file specifying every varying quantity, seed count, budget, and metric threshold); the researcher runs the
experiment on hardware, recording raw results and resolved config alongside git state; an adversarial reviewer (``Tarka'')
recomputes every number from the raw results files, checks match-for-match fairness in parameters, tokens, optimizer,
and learning rate, and attacks the mechanism interpretation. If the reviewer finds a material defect (confounding,
measurement error, or mechanistic misread), the claim is either refuted, rescoped, or re-run. Only after this review
closes may a finding enter the paper. Two reviews materially overturned author conclusions (Sections~\ref{sec:scale} and
\ref{sec:closedloop}), demonstrating the audit's teeth. Hypotheses carry named outcomes: \textsc{met}, \textsc{partial},
or \textsc{missed}. A well-documented miss is a valid closure state and often contains the most design information.

\subsection{Hypotheses and Verdicts}

Table~\ref{tab:hypotheses} states the six pre-registered hypotheses (H1–H6) with their verbatim quantitative gates,
the experiment verdicts, and the findings that support them. H7–H10, registered for the open-weight 1–3B scale program,
are marked in progress.

\begin{table}[H]
\centering
\small
\caption{Pre-registered hypotheses with quantitative gates, verdicts, and finding references. Gates were locked
before any experiment ran; verdicts are named outcomes (MET/PARTIAL/MISSED) assigned by adversarial review. H7–H10
target 1–3B-parameter models and are in progress.}
\label{tab:hypotheses}
\begin{tabular}{@{}p{1.0cm}p{3.8cm}p{2.6cm}p{1.4cm}@{}}
\toprule
\textbf{ID} & \textbf{Gate (quantitative threshold)} & \textbf{Outcome} & \textbf{Findings} \\
\midrule
H1 & EqLM $\geq$95\% of explicit BLiMP at matched params/tokens (121M) & \textsc{missed} & F13–F20 \\
H2 & MMD last-iterate convergence where GDA cycles: $R^2 \geq 0.9$ on matrix games & \textsc{met} & F1–F3 \\
H3 & PMA $\geq$ DPO held-out accuracy AND lower KL-to-reference drift & \textsc{partial} & F21 \\
H4 & Auction beats best single specialist on mixed-domain perplexity & \textsc{met} & F22 \\
H5 & Auction advantage persists in closed-loop greedy generation & \textsc{missed} & F23 \\
H6 & Recover $\geq$50\% of width-gap + certification ($R > 0.5$) at $\leq$12 iterations & \textsc{partial} & F24 \\
H7–H10 & 1–3B conversion and preference-tuning (in progress) & \textsc{in progress} & — \\
\bottomrule
\end{tabular}
\end{table}

\subsection{Corpora}

The pretraining program uses the BabyLM strict-small dataset (10M unique tokens, drawn from the BabyLM-2026
benchmark). All arms—explicit baseline, EqLM implicit, EqLM unrolled, and ablation variants—train on the identical
token stream, with every experiment fixing the sequence length to 128 tokens and the batch size to 32. This ensures
matched token budgets across model variants. At 121M parameters, the full pretraining run (10,000 steps per arm)
consumes approximately 40.96M token presentations ($10{,}000 \times 32 \times 128$). Evaluation harnesses (BLiMP
and lm-evaluation-harness) use held-out streams disjoint from training.

For the 1–3B scale program (H7–H10), the FineWeb-Edu corpus is used, selected for compatibility with published
open-weight baseline checkpoints (e.g.\ Qwen3-1.7B) that will serve as the conversion source. This guarantees
that both the pretrained base and the EqLM conversion run on identically sourced data, removing corpus as a
confound in downstream evaluation.

\subsection{Evaluation Harness and Benchmark}

At 11M–121M parameters, the core evaluation metric is BLiMP (Basic Language Corpus Minimal Pairs), a test set of
1,335 grammatical acceptability pairs spanning agreement, control, binding, tense, and relative clause phenomena.
All BLiMP scores reported in this paper are computed over a held-out 1,000-pair subset, held constant across seeds.
The minimal-pair format provides a binomial resolution: at $n = 1{,}000$ pairs, the standard error of the proportion
is approximately $\sqrt{p(1-p)/n} \approx 0.014$ (1.4 percentage points at one standard deviation for $p \approx 0.5$).
Confidence intervals use bootstrap resampling over the 1,000 pairs, stratified by seed.

For the open-weight 1–3B scale program, evaluation transitions to the lm-evaluation-harness framework. On the
user's hardware (Qwen3-1.7B baseline on a GB10 node), the zero-shot base rates measured for the harness's canonical
benchmarks are: ARC-Challenge accuracy 0.4067, ARC-Challenge acc\_norm 0.4433; HellaSwag accuracy 0.43, HellaSwag
acc\_norm 0.5067; GSM8K flexible-extract 0.4567. These base rates are held constant as reference points; all
comparisons—EqLM conversion vs.\ explicit baseline, preference-optimized variants, auction aggregates—use the same
harness and the same evaluation hardware to prevent measurement confounds that arise from different systems or code
paths. Rather than lifting numbers from published model cards, the program re-measures every baseline on the project's
own hardware under the same code path (kinetic\_ai/serve/), ensuring that architectural comparisons are not confounded
by harness implementation, prompt format, or hardware differences.

\subsection{Hardware and Thermal Constraints}

Training of the 121M-parameter models (exp10, exp13) runs on an NVIDIA RTX 5090 GPU with 32 GB of VRAM. The machine
hosting this GPU (a dedicated DGX Spark node) is reserved exclusively for this program's training jobs. Evaluation and
inference—including BLiMP scoring, checkpoint conversion, and trace serving for the application—runs on a GB10 node
with 121 GB of unified GPU memory. The GB10 operates under a software thermal governor triggered by a hardware cooling
defect that causes the platform to shut down around 91°C. This constraint is not a limitation of the architecture but
a hardware-level fact of the test environment; it is documented because memory and wall-time measurements reflect the
thermal throttling. Peak activation memory is reported as the deterministic allocator peak for a fixed computation
graph, an architecture-level measurement rather than a per-run distribution. All memory figures are verified bit-identical
across seeds, consistent with the deterministic-allocator design.

\subsection{Statistical Protocol}

Every verdict in Table~\ref{tab:hypotheses} is grounded in one of three statistical frameworks, depending on the
research question.

\paragraph{Paired bootstrap over seeds (main pretraining and preference arms).} When two architectures or training
regimes are evaluated on identical data across multiple random seeds, the per-seed measurements (e.g.\ BLiMP scores
0.682, 0.675, 0.693 for three seeds) form the basis. The bootstrap confidence interval is computed by resampling the
seed-level means 10,000 times with replacement and recording the 2.5th and 97.5th percentiles. For pairwise comparisons
(e.g.\ H1's verdict that EqLM reaches 0.930 of baseline), the per-seed ratios are bootstrapped. A one-sided paired
$t$-test (where appropriate for a directional claim) uses the seed-level data; the sign test (an exact paired
permutation test) is reported as well when $n$ is small (e.g.\ $n = 5$ prompts).

\paragraph{Binomial test (steering and causal interventions).} When testing whether an intervention produces an effect
at above-chance rates (e.g.\ H2b: does steering a circuit-selected feature flip the top-1 prediction more often than
a random feature?), the data are counts of successes and failures across a fixed set of trials. A one-sided binomial
test is used with a null hypothesis of random performance (e.g.\ 1\% per-token chance or 50\% for a two-way choice).
The null-hypothesis probability is specified pre-registered in the hypothesis gate.

\paragraph{Like-for-like audit (all comparisons).} After the experiment runs and before any verdict is assigned, the
adversarial reviewer checks that all arms being compared satisfy strict match criteria: identical data stream (same
tokenized, shuffled corpus file), identical steps and batch size, identical optimizer and learning rate, identical
gradient clipping and warmup schedules, and parameter count within 5\%. The audit also recomputes key aggregate
quantities (means, confidence intervals) directly from the raw per-seed result files. Where the reviewer detects a
confound—different data loaders, different learning rates, different tokenization—the claim is either rescoped or
the arms are re-run under matched conditions. Three findings underwent major rescoping after audit (F15, F20, F24);
two were refuted and rewritten (F14's "no fixed point" discovery, F20's mechanism claim); and one closed-loop
experiment was run entirely because F22's advertised gate on "closed-loop auction" would have failed, leading the
reviewer to demand the experiment and rescope F22 itself.

\subsection{Configuration Hashing and Provenance}

Every experiment run produces a results JSON file containing (i) per-seed raw measurements (loss, BLiMP scores, wall
time, peak memory); (ii) a resolved configuration dictionary (all hyperparameters, seeds, batch size, corpus, and
paths); (iii) a SHA-256 hash of the resolved configuration; and (iv) the git commit hash at run time. The findings
ledger (research/memory/findings.md) records, for each finding, the experiment name, the config hash, the commit,
the data range (which seed IDs, which steps, which pairs), and the extracted number. This ensures that every reported
quantity is reproducible and that its context (hardware, code version, data version) is explicit. If an experiment is
re-run (e.g.\ due to a Tarka-identified confound), both the old and new config hashes are recorded, and the finding's
entry notes which hash applies to the reported number.

\subsection{Adversarial Review Procedure}

The adversarial reviewer operates under a standing mandate: recompute every number from the raw result files; flag any
arithmetic or measurement error; challenge the mechanistic interpretation; and demand re-analysis or re-running if the
evidence does not support the claim. The reviewer may refute a claim, rescope it, demand additional experiments (e.g.\
per-domain breakdowns, ablations), or approve it. Review verdicts are recorded and published alongside the finding.
Two high-impact examples appear in this paper. First, F20's mechanistic claim—``solver convergence fails at 121M
width''—was refuted by the reviewer, who showed that solver convergence rate was 0.0 at both 11M and 121M, so the
failure was not binary but graded: the map's contraction per iteration decreases with width under a fixed 12-iteration
budget, widening the truncation penalty from 7\% to 21\%. Second, F22's gate stated that "closed-loop generation"
achieves the auction advantage; the reviewer demanded that this gate be re-scoped to "scoring-time selection" based
on the teacher-forced experimental design, and then insisted that the pre-registered closed-loop experiment (H5) be
run. H5 was missed, confirming the reviewer's skepticism and validating F22's rescoped claim.

\subsection{Scale and Token Limits}

All results reported in Sections~\ref{sec:convergence} through~\ref{sec:auction} operate at 11M–124M parameters on
10M words of training data. The 1–3B scale program (H7–H10) uses FineWeb-Edu and is in progress; uptraining timelines
and conversion-damage sweeps have been completed (F25), but training and preference-tuning runs are ongoing. Results
are reported only when they have closed with adversarial review and operator sign-off.
